\documentclass[11pt]{ctexart}
\usepackage{amsmath,amssymb,longtable}
\begin{document}
\section*{P161: 当前完整数学模型（2026-10-20）}
\section*{修订依据与边界}按源文第14–15页及物理服务定义补齐直接/同机两段行程、完整电量与有限停机坪/航路容量。时槽上限来自时域和最短飞行+周转的证明；请求集合由实例完整给定。

\section*{数据、时间与飞行}
节点$V=V_P\cup V_S$、飞机$N$、乘客组$A$、基地$b_n$、飞行时间$r_{ij}$、
能耗$e_{ij}$、成本$c_{ij}$、时域$H$、周转$t_e$、座位、停机坪和航路容量均为实例参数。
使用实例精确分钟刻度（当前0.1分钟），飞行时长和时域均为整数刻度。
全部有效飞行槽上限
$S=\min\{S^{input},\lfloor H/(t_e+\min r_{ij})\rfloor\}$来自每段消耗至少周转和最短飞行时间。
二元$X_{ijmnt}$表示飞机$n$的第$m$段在$t$起飞；
$x_{ijmn}=\sum_tX_{ijmnt}$，$h_{mn}=\sum_{ij}x_{ijmn}$。
只保留$i\ne j$、$mt_e+(m-1)\min r\le t\le H-r_{ij}$的起飞；
离开非基地后须有最少返回时间；第一段从基地，最后有效槽到基地。
\[
h_{1n}=1,\quad h_{mn}\le1,\quad h_{mn}\ge h_{m+1,n},\quad
\sum_ix_{ijmn}\ge\sum_ix_{ji,m+1,n},
\quad\sum_{m,j}x_{b_njmn}=\sum_{m,i}x_{ib_nmn}.
\]
虚拟槽0的路线和乘客量为零，$a_{0n}=d_{0n}=0$。
活跃槽$a_{mn}=\sum_{ijt}(t+r_{ij})X_{ijmnt}$；
非活跃槽$a_{mn}=a_{m-1,n}$。所有槽
$a_{mn}\ge a_{m-1,n}+t_eh_{mn}+\sum_{ij}r_{ij}x_{ijmn}$，
$d_{mn}=a_{mn}-\sum_{ij}r_{ij}x_{ijmn}$，$a,d\in[0,H]$。
\section*{电池与充电}
$u_{0n}=\bar E$，$u_{mn}\in[\underline E,\bar E]$，$q_{mn}\ge0$，
\[
u_{mn}=u_{m-1,n}+q_{mn}-\sum_{ij}e_{ij}x_{ijmn},\quad
u_{m-1,n}+q_{mn}\le\bar E,\quad q_{mn}\le\bar Eh_{mn}.
\]
$q_{1n}=0$；后续$q_{mn}\le\bar E\sum_{i\in V_P,j}x_{ijmn}$，
$q_{mn}\le rate\cdot tick\cdot(d_{mn}-a_{m-1,n})$。
\section*{完整旅客行程}
每组$a$仅能选择一个直接$(o_a,d_a,m,n)$或允许的两段
$(o_a,h,m,n),(h,d_a,m+1,n)$（$h$不同于起终点），对应二元$z$小于每段$x$。
首段前次到达$a_{m-1,n}\le request_a$；首发在
$[request_a,\min(latest_a,request_a+w)]$，两段同机中转间隔恰等$t_e$。
每个飞行槽最多2组，所载人数不超过座位容量；每组总选择$\sum z\le1$。
\section*{真实时空占用}
飞行$X$占用全部整数刻度$t+1,\ldots,t+r_{ij}$；
每一有向航路每刻度的占用数不超过有限$cap_{flt}$。
飞机在节点的停放量为初始基地指示，加所有到达时间严格小于当前刻度的入飞，
减所有出发时间严格小于当前刻度的出飞；其值在$[0,1]$。
每节点每刻度停放总量不超过给定节点容量。这覆盖整个时域和所有转场，不仅起飞计数。
\[
\max\sum_a[(fare_a+penalty_a)\sum z_a-penalty_a]
-\sum_{ijmn}c_{ij}x_{ijmn}.
\]
$fare_a$按输入的可接受中转类型取相应票价，即使该允许中转组最终乘直达也按该报价；
所有原请求均有服务或拒绝的明确记录。

\end{document}
